\documentclass[10pt,a4paper]{amsart}
\usepackage[margin=19mm]{geometry}
\usepackage{amsmath,amssymb,mathtools}
\usepackage[scr=rsfs,cal=euler]{mathalfa}
\usepackage{xfrac}
\usepackage[unicode,hidelinks,bookmarks=false]{hyperref}
\newcommand{\ie}{i.e.\ }
\newcommand{\cf}{cf.\ }
\newcommand{\resp}{respectively\ }
\newcommand{\Subsection}{\subsection}
\providecommand{\nobreakdash}{\nobreak\hbox{-}}
\setlength{\emergencystretch}{2em}
\multlinegap0pt
\usepackage{bm}
\usepackage{bbm}
\usepackage{dsfont}
\usepackage{xspace}
\usepackage{cancel}
\usepackage{mathtools}
\usepackage{amsthm}
\makeatletter
\@ifundefined{theorem}{\newtheorem{theorem}{Theorem}}{}
\@ifundefined{prop}{\newtheorem{prop}[theorem]{Proposition}}{}
\makeatother
\newcommand{\Inn}{\operatorname{Inn}}
\newcommand{\Z}{\operatorname{Z}}
\title[Fundamental $n$-quandles of links are residually finite]{Fundamental $n$-quandles of links are residually finite}
\author{Neeraj Kumar Dhanwani, Deepanshi Saraf,, Mahender Singh}
\date{Source: arXiv:2403.17703v2 (CC BY 4.0)}
\begin{document}
\maketitle

\noindent\emph{Source and licence.} Fundamental $n$-quandles of links are residually finite. Version arXiv:2403.17703v2 (CC BY 4.0), distributed under the Creative Commons Attribution 4.0 International licence, as recorded in the arXiv licence metadata for this exact version and shown on the abstract page https://arxiv.org/abs/2403.17703v2. Full rights notice: licenses/arxiv-2403.17703-CC-BY-4.0.txt.\par\smallskip

\noindent\textit{Editorial scope.} Counted unit: the Prop whose source label is \texttt{twisted union subquandle sep} in the article (source0.tex lines 564--566, source label \texttt{twisted union subquandle sep}), delivered together with the article's own complete original proof of it (source0.tex lines 568--586, one \texttt{proof} environment of 19 lines). No mathematics is written, repaired or re-derived: statement and proof are the article's own text line for line. Required context retained uncounted: none. Every definition, display and label of those lines is retained unchanged. Omitted from this package and not redistributed: 1--563, 567--567, 587--717. Nothing inside a retained excerpt is omitted or altered: every retained body is delivered exactly as the article's lines are written, so no omission receipt of any kind accompanies this package. Citations: the retained excerpts carry no citation command at all, so no bibliography entry of the article is transcribed and no reference is redistributed. Reference adaptation: none was needed and none was made; every reference inside the delivered unit resolves inside it, so no unresolved reference or citation is produced. Printed slips kept literal: no printed word, symbol, hypothesis or number of the counted statement or of its proof was altered. Layout and class: the article's own class is \texttt{amsart} with packages including amscd, amsfonts, amsmath, amssymb, amsthm, enumerate, flexisym, float, graphicx, hyperref, pdfpages, tikz, mathtools, comment; the wrapper is the lane's pinned \texttt{amsart} editorial wrapper at 10pt on A4, which restates the theorem-like environments the retained text needs and adds the article's own macro definitions that the retained text needs (\texttt{\textbackslash Inn}, \texttt{\textbackslash Z}), with extra wrapper packages (bm, bbm, dsfont, xspace, cancel, mathtools). The delivered numbering is the unit's own. Assets: no retained span contains a figure environment, a graphics-inclusion command, a PDF-page inclusion, a table or a TikZ picture; no image file is copied into this package, the package declares no asset dependency and no image file is redistributed with it. Licence: version arXiv:2403.17703v2 of this article is distributed under the Creative Commons Attribution 4.0 International licence, as recorded in the arXiv licence metadata for this exact version and shown on the abstract page https://arxiv.org/abs/2403.17703v2. A full rights notice is delivered at licenses/arxiv-2403.17703-CC-BY-4.0.txt. Characters: every retained excerpt is pure ASCII, so the delivered engine, XeLaTeX with its default Latin Modern text font, reports no missing character.
\begin{prop}\label{twisted union subquandle sep}
Let $(X_1,\star_1)$ and $(X_2,\star_2)$ be subquandle separable quandles. If $f\in \Z(\Inn(X_1))$ and $g\in \Z(\Inn(X_2))$ are finite order elements, then $(X,\ast)$ is subquandle separable.
\end{prop}

\begin{proof}
Let $Y$ be a  finitely generated subquandle of $X$ and $z\in X\setminus Y$. Without loss of generality, we can assume that $z\in X_1$. We first claim that $Y\cap X_1$ is a finitely generated subquandle of $X_1$. Suppose that $Y$ is generated by $\{a_1,a_2,\ldots,a_k,b_1,b_2,\ldots,b_\ell\}$, where each $a_i\in X_1$ and each $b_j\in X_2$. Suppose that the order of $f$ is $m$. Since $f\in \Z(\Inn(X_1))$, a direct check using the quandle operation in $X$ shows that $Y\cap X_1$ is generated by the finite set $\{f^{t_i}(a_i)\mid 1\leq i\leq k, ~ 1\leq t_i\leq m\}$.  Since $z\in X_1$ and $Y\cap X_1$ is a finitely generated separable subquandle of $X_1$,  there is a surjective quandle homomorphism $\phi:X_1\to F$, where $(F,\circ)$ is a finite quandle, such that $\phi(z)\notin \phi(Y\cap X_1)$. Let $p$ be a symbol disjoint from $F$ and $F'=F\sqcup\{p\}$. Define a binary operation $\circ'$ on $F'$ as 
$$x\circ' y=\begin{cases}
			x\circ y \quad\text{ if } x,y\in F,\\
			p   \quad\quad\quad\text{ if } x=p,\\
			\phi(f(z)) \quad\text{ if } y=p,~x \ne p
		\end{cases}$$
where $z \in X_1$ is such that $\phi(z)=x$. We claim that $(F',\circ')$ is a quandle. Suppose that $z_1, z_2 \in X_1$ such that $\phi(z_1)=\phi(z_2)=x$. Since $f \in \Inn(X_1)$,
we see that  $\phi(f(z_1))= \phi(f(z_2))$, and the binary operation $\circ'$ is indeed well-defined. Let $S_p: F' \to F'$ be the right multiplication by $p$. For arbitrary $x_1, x_2 \in F$, let $z_1, z_2 \in X_1$ such that $\phi(z_1)=x_1$ and  $\phi(z_2)=x_2$.  Suppose that $S_p(x_1)=S_p(x_2)$, that is, $x_1 \circ' p= x_2 \circ' p$. This gives $\phi(f(z_1))= \phi(f(z_2))$. Since $f \in \Inn(X_1)$, it follows that $x_1=\phi(z_1)=\phi(z_2)=x_2$. Thus, $S_p$ is injective, and finiteness of $F'$ implies that it is a bijection of $F'$. Further, we see that 
$$S_p(x_1 \circ' x_2)= S_p(x_1 \circ x_2)= \phi f (z_1 \star_1 z_2) =\phi f(z_1) \circ' \phi f(z_2) = S_p(x_1) \circ' S_p(x_2),$$
$$S_p(x_1 \circ' p)= S_p(\phi f(z_1))= \phi f(f(z_1)) =\phi f(z_1) \circ' p = S_p(x_1) \circ' S_p(p)$$
and
$$S_p(p \circ' x_2)= S_p(p)=p =p \circ' S_p(x_2) = S_p(p) \circ' S_p(x_2).$$
Using the fact that $f \in \Z(\Inn(X_1))$, we have
$$S_{x_1}(x_2 \circ' p)= S_{x_1}(\phi f (z_2))=\phi S_{z_1} f(z_2)=\phi f S_{z_1}(z_2) = S_{x_1}(x_2) \circ' p = S_{x_1}(x_2) \circ' S_{x_1}(p)$$
and
$$S_{x_1}(p \circ' x_2)= S_{x_1}(p)=p= S_{x_1}(p) \circ' S_{x_1}(x_2).$$
This proves our claim. It is easy to see that the map $\Phi:X\to F'$ defined by $\Phi(X_2)=p$ and $\Phi(y)=\phi(y)$ for $y\in X_1$,  is a quandle homomorphism. Further, $\Phi(z)\notin \Phi(Y)$, and the result follows.
\end{proof}

\end{document}
